% ============================================================
% CHAPTER 5 — Project 2: A Complete CLI Application
% ============================================================
\chapter{Project 2: A Complete CLI Application}\index{CLI}

\section*{What You Will Build}
A \textbf{command-line Task Manager} --- a professional Python application that:
\begin{itemize}
  \item Manages tasks with CRUD operations (create, read, update, delete)
  \item Persists data to JSON files
  \item Supports filters and search
  \item Has a priority and deadline system
  \item Includes automated tests with full coverage
\end{itemize}

\textbf{Estimated time}: 30--45 minutes

% ---------------------------------------------------------
\section{The Professional Context (ADLC Phases 0--2)}

This project is the first one where \texttt{\_CONTEXT.md} really makes the difference between a mediocre result and a professional one.

\begin{praticabox}[title={Hands-on --- Project setup}]
\begin{enumerate}
  \item Create the folder \texttt{taskmaster}
  \item Open it in VS~Code
  \item Create the file \texttt{\_CONTEXT.md}:
\end{enumerate}
\end{praticabox}

\begin{lstlisting}[language={},title={\texttt{\_CONTEXT.md} for TaskMaster CLI}]
# Progetto: TaskMaster CLI

## Scopo
Applicazione CLI per la gestione di task personali con
persistenza locale. Strumento per uso quotidiano dal terminale.

## Tecnologie
- Python 3.11+
- SOLO standard library (argparse, json, os, datetime, pathlib, enum)
- Testing: pytest
- Nessuna dipendenza esterna ammessa

## Struttura del Progetto
taskmaster/
+-- _CONTEXT.md
+-- main.py              <- Entry point (solo parsing argomenti)
+-- task_manager.py      <- Classe TaskManager: logica CRUD
+-- models.py            <- Dataclass Task con tutti i campi
+-- storage.py           <- Classe Storage: lettura/scrittura JSON
+-- formatters.py        <- Formattazione output per il terminale
+-- tests/
|   +-- conftest.py      <- Fixtures condivise
|   +-- test_task_manager.py
|   +-- test_storage.py
|   +-- test_formatters.py
+-- data/
    +-- tasks.json       <- Creato automaticamente al primo uso

## Modello Dati (Task)
- id: int (auto-incrementale, non riutilizzato)
- title: str (obbligatorio, max 100 caratteri)
- description: str (opzionale, default "")
- status: enum [todo, in_progress, done] (default: todo)
- priority: enum [low, medium, high] (default: medium)
- created_at: str ISO 8601
- updated_at: str ISO 8601
- due_date: str ISO 8601 o null (opzionale)

## Comandi CLI
taskmaster add "Titolo del task"
taskmaster add "Titolo" -d "Descrizione" -p high
taskmaster add "Titolo" --due 2026-04-15
taskmaster list
taskmaster list --status todo
taskmaster list --priority high
taskmaster list --overdue
taskmaster show 3
taskmaster update 3 --status done
taskmaster delete 3
taskmaster stats

## Convenzioni di Codice
- Naming: snake_case per tutto
- Classi: PascalCase (TaskManager, Storage, Task)
- Type hints: OBBLIGATORI su tutte le firme
- Docstring: su ogni classe e metodo pubblico
- Enum: usa enum.Enum per status e priority
- Dataclass: usa @dataclass per il modello Task
- Pathlib: usa pathlib.Path, mai os.path

## Vincoli di Implementazione
- NON usare variabili globali.
- NON ignorare le eccezioni.
- NON usare exit() nel mezzo del codice. Solo in main.py.
- Il file tasks.json DEVE essere creato automaticamente.
- L'ID dei task DEVE essere auto-incrementale e MAI riutilizzato.
- Il formato JSON deve essere indentato (indent=2).
- La cancellazione di un task chiede conferma (y/N).
\end{lstlisting}

\begin{notebox}
Notice how this \texttt{\_CONTEXT.md} is much more detailed than the one in Chapter~4. It defines the data model, the exact CLI commands, the code conventions, and the implementation constraints. The more detailed the context is, the better the generated code will be.
\end{notebox}

% ---------------------------------------------------------
\section{Code Generation (ADLC Phases 3--4)}

\begin{praticabox}[title={Hands-on --- Phase 3: Proof of Value}]
Before generating the entire project, do a quick test. In Copilot Agent Mode:
\begin{lstlisting}[language={}]
Leggi il _CONTEXT.md. Come primo passo, genera solo il file
models.py con la dataclass Task e gli enum Status e Priority.
Non generare nient'altro per ora.
\end{lstlisting}
Verify that it uses \texttt{@dataclass}, \texttt{enum.Enum}, type hints, and default values.
\end{praticabox}

\begin{praticabox}[title={Hands-on --- Phase 4: Full generation}]
\begin{lstlisting}[language={}]
Perfetto. Ora implementa l'intero progetto secondo il
_CONTEXT.md. Procedi in quest'ordine:
1. storage.py (persistenza JSON)
2. task_manager.py (logica CRUD)
3. formatters.py (output terminale)
4. main.py (CLI con argparse)
5. tests/ (tutti i file di test)
\end{lstlisting}
\end{praticabox}

\begin{tipbox}
Specifying the implementation order helps the AI build the project layer by layer, avoiding circular references and missing dependencies.
\end{tipbox}

% ---------------------------------------------------------
\section{Reviewing the Generated Code}

After Copilot has generated all the files, do a systematic review.

\subsection*{Review checklist}

\textbf{\texttt{models.py}}:
\begin{itemize}
  \item \texttt{Status} and \texttt{Priority} are Enums, not strings
  \item \texttt{Task} is a dataclass with all the fields from the context
  \item Methods \texttt{to\_dict()} and \texttt{from\_dict()} for serialization
  \item \texttt{created\_at} and \texttt{updated\_at} are generated automatically
\end{itemize}

\textbf{\texttt{storage.py}}:
\begin{itemize}
  \item Uses \texttt{pathlib.Path}, not \texttt{os.path}
  \item Creates the JSON file automatically if it does not exist
  \item Reads and writes with \texttt{indent=2}
  \item Handles \texttt{FileNotFoundError} and \texttt{JSONDecodeError}
\end{itemize}

\textbf{\texttt{task\_manager.py}}:
\begin{itemize}
  \item The ID is auto-incremental and is never reused
  \item Supports all commands from the context
  \item \texttt{update\_task} updates \texttt{updated\_at}
\end{itemize}

\textbf{\texttt{main.py}}:
\begin{itemize}
  \item Uses \texttt{argparse} with subcommands
  \item Helpful messages, not stack traces
\end{itemize}

\textbf{\texttt{tests/}}:
\begin{itemize}
  \item Uses \texttt{tmp\_path} for storage
  \item Tests positive cases and edge cases
  \item Each test is independent
\end{itemize}

\begin{warnbox}
If you find a problem, do not edit the code by hand. Describe the issue to Copilot and ask it to fix it. This is the 0-code workflow.
\end{warnbox}

% ---------------------------------------------------------
\section{Testing and Execution}

\begin{praticabox}[title={Hands-on --- Testing}]
\begin{lstlisting}[language=bash]
python -m pytest tests/ -v
\end{lstlisting}
If any test fails, copy the error into Copilot Chat:
\begin{lstlisting}[language={}]
Questo test fallisce con il seguente errore:
[incolla l'errore]
Correggi il codice per far passare il test.
\end{lstlisting}
\end{praticabox}

\begin{praticabox}[title={Hands-on --- Using the application}]
\begin{lstlisting}[language=bash]
# Aggiungi task
python main.py add "Leggere il capitolo 6" -p high
python main.py add "Comprare il latte" --due 2026-04-01

# Lista tutti
python main.py list

# Filtra
python main.py list --status todo
python main.py list --priority high

# Aggiorna e completa
python main.py update 1 --status in_progress
python main.py update 1 --status done

# Statistiche
python main.py stats
\end{lstlisting}
\end{praticabox}

\begin{checkpointbox}
If all commands work and the tests pass, you have a complete professional CLI application.
\end{checkpointbox}

% ---------------------------------------------------------
\section{Iteration: Adding Features}

\begin{praticabox}[title={Hands-on --- Export to Markdown format}]
\begin{lstlisting}[language={}]
Aggiungi un comando "export" che esporta tutti i task in un
file Markdown formattato come tabella.
Uso: python main.py export report.md
Il file deve contenere una tabella con i task organizzati
per priorita' (high -> medium -> low).
Aggiungi i test corrispondenti.
\end{lstlisting}
\end{praticabox}

\begin{praticabox}[title={Hands-on --- Text search}]
\begin{lstlisting}[language={}]
Aggiungi un comando "search" che cerca tra titoli e
descrizioni. Uso: python main.py search "email"
La ricerca deve essere case-insensitive.
Aggiungi i test corrispondenti.
\end{lstlisting}
\end{praticabox}

% ---------------------------------------------------------
\section{The Pattern You Will Use Forever}

This project established the pattern you will follow throughout the book:

\begin{enumerate}
  \item \textbf{\texttt{\_CONTEXT.md}} $\rightarrow$ Define the project (5--15~min)
  \item \textbf{Proof of Value} $\rightarrow$ Generate a small piece and verify (5~min)
  \item \textbf{Generation} $\rightarrow$ The AI implements everything (10--20~min)
  \item \textbf{Review} $\rightarrow$ You check the quality (5--10~min)
  \item \textbf{Testing} $\rightarrow$ Verify that it works (2--5~min)
  \item \textbf{Iteration} $\rightarrow$ Add features and fix issues (ongoing)
  \item \textbf{Update} $\rightarrow$ Improve the \texttt{\_CONTEXT.md} (2~min)
\end{enumerate}

This pattern scales: it works for a 200-line CLI tool just as well as for a 10,000-line full-stack app.

% ---------------------------------------------------------
\section{Lessons Learned (ADLC Phase 6)}

\begin{lstlisting}[language={},title={Lessons to add to the context}]
## Lezioni Apprese
- Specificare l'ordine di implementazione
  (storage -> logic -> UI -> main -> tests)
- I test con tmp_path devono creare un'istanza fresh
- argparse con subparsers richiede set_defaults(func=...)
- L'IA tende a dimenticare la conferma y/N sulla
  cancellazione: metterla nei vincoli
\end{lstlisting}

% ---------------------------------------------------------
\section*{Project Summary}

\begin{center}
\begin{tabularx}{\textwidth}{l X}
\toprule
\textbf{Aspect} & \textbf{Detail} \\
\midrule
Language & Python 3.11+ \\
Dependencies & Zero (standard library only) \\
Generated files & $\sim$7 files, $\sim$400--600 total lines \\
Tests & $\sim$15--20 test cases \\
Total time & $\sim$30--45 minutes \\
Code written by hand & Only \texttt{\_CONTEXT.md} \\
\bottomrule
\end{tabularx}
\end{center}
